\begin{helper}
Fix one pair of support labels \(B,J\) and the terminal coordinate set
\(U\).  Let \(\mathcal P\) be the transcript family.  For a realized
branch/transcript pair \((\beta,\tau)\), write
\[
P_\tau,\qquad A_\tau,\qquad Z_\tau,\qquad F_{\beta,\tau}
\]
for the forced-present set, forced-absent set, selected absent set, and
fixed present support.  Assume the fixed realized geometry
\[
F_{\beta,\tau}\subseteq P_\tau,\qquad
P_\tau,A_\tau\subseteq U,\qquad
P_\tau\cap A_\tau=\emptyset,\qquad
Z_\tau\subseteq A_\tau .
\]
Assume also the prefix-scan interface used by
\(\noderef{FixedSetHistoryCellRedPrefixScanInvariant}\): each test at a
coordinate depends only on the earlier prefix, compatible complete
assignments determine their blue branch and transcript unless a
selected-present sibling is active, selected-present siblings lie in
\(Z_\tau\), and an active selected-present sibling in a realized compatible
cell rules out compatibility with every realized cell.

Finally assume the first-mismatch certificate for the fixed residual
scan.  Namely, if one residual assignment \(A_{\rm res}\subseteq U\) lies
in the residual cylinders of two realized pairs \((\beta,\tau)\) and
\((\beta',\tau')\), the two reconstructed complete assignments
\[
A^\ast=(A_{\rm res}\setminus Z_\tau)\cup F_{\beta,\tau},
\qquad
A^{\ast\prime}=(A_{\rm res}\setminus Z_{\tau'})
  \cup F_{\beta',\tau'}
\]
are compatible with their respective realized cells, and
\((\beta,\tau)\ne(\beta',\tau')\), then a selected-present sibling is
active in at least one of the two reconstructed compatible assignments.

Under these hypotheses, the same residual assignment has a unique realized
owner.  More explicitly, suppose
\[
P_\tau\setminus F_{\beta,\tau}\subseteq A_{\rm res},\qquad
A_{\rm res}\cap(A_\tau\setminus Z_\tau)=\emptyset ,
\]
and
\[
P_{\tau'}\setminus F_{\beta',\tau'}\subseteq A_{\rm res},\qquad
A_{\rm res}\cap(A_{\tau'}\setminus Z_{\tau'})=\emptyset .
\]
If \(A^\ast\) is compatible with \((\beta,\tau)\) and
\(A^{\ast\prime}\) is compatible with \((\beta',\tau')\), then
\(\beta=\beta'\) and \(\tau=\tau'\).
\end{helper}
\begin{proof}
Apply \(\noderef{FixedSetHistoryCellRedPrefixScanInvariant}\) to the fixed
prefix-scan interface.  Its first conclusion is the pruning clause: if a
selected-present sibling is active in a realized compatible complete
assignment, then no realized cell is compatible with that same complete
assignment.  Its second conclusion is the branch/transcript determinism for
compatible complete assignments.  Only the pruning clause is needed for
this node.

Assume for contradiction that the two owners are different, so either
\(\beta\ne\beta'\) or \(\tau\ne\tau'\).  The first-mismatch certificate
applies to the two residual-cylinder memberships and to the two
compatibility certificates.  It produces an active selected-present sibling
in one of the two reconstructed assignments.

If the sibling occurs in \(A^\ast\), the pruning clause applied to the
realized cell \((\beta,\tau)\) says that \(A^\ast\) is compatible with no
realized cell, contradicting the assumed compatibility of \(A^\ast\) with
\((\beta,\tau)\).  If the sibling occurs in \(A^{\ast\prime}\), the same
argument contradicts compatibility of \(A^{\ast\prime}\) with
\((\beta',\tau')\).  Both alternatives are impossible.  Hence the assumed
owner mismatch is impossible, and therefore \(\beta=\beta'\) and
\(\tau=\tau'\).
\end{proof}
